% ------------------------------------------
% TRANG BÌA (Giống hệt file mẫu PDF)
% ------------------------------------------
\begin{titlepage}
    \begin{center}
        {\large ĐẠI HỌC BÁCH KHOA HÀ NỘI} \\
        \vspace{1cm}
        \textbf{\Large ĐỒ ÁN TỐT NGHIỆP} \\
        \vspace{2cm}
        {\LARGE \textbf{Xây dựng hệ thống chấm điểm phiếu trắc nghiệm bằng hình ảnh}} \\
        \vspace{2cm}
        \textbf{\large Phạm Quốc Dũng} \\
        Dung.pq225817@sis.hust.edu.vn \\
        \vspace{1cm}
        \textbf{Ngành Công nghệ thông tin Việt Nhật} \\
        \textbf{Trường:} Công nghệ Thông tin và Truyền thông
    \end{center}

   \vspace{1.5cm}
\begin{flushright}
    \begin{minipage}{6cm}
        \centering
        \textbf{Giảng viên hướng dẫn} \\
        TS. Trần Việt Trung \\
        \vspace{1.2cm}
        Chữ kí GVHD
    \end{minipage}
\end{flushright}

    \vfill
    \begin{center}
        HÀ NỘI, 06/2025
    \end{center}
\end{titlepage}

% ------------------------------------------
% PHẦN THỦ TỤC
% ------------------------------------------
\pagenumbering{roman}

%

% 2. Lời cam kết
\chapter*{LỜI CAM KẾT}
\addcontentsline{toc}{chapter}{LỜI CAM KẾT}
\begin{flushleft}
    \textbf{Họ và tên sinh viên:} Phạm Quốc Dũng \\
    \textbf{Điện thoại liên lạc:} 0343624070 \\
    \textbf{Email:} Dung.pq225817@sis.hust.edu.vn \\
    \textbf{Lớp:} Công nghệ thông tin Việt Nhật 04 - K67 \\
    \textbf{Hệ đào tạo:} Đại học chính quy
\end{flushleft}
\vspace{0.5cm}
Tôi – Phạm Quốc Dũng – xin cam kết rằng Đồ án Tốt nghiệp (ĐATN) là công trình nghiên cứu do chính tôi thực hiện dưới sự hướng dẫn của TS. Trần Việt Trung. Các kết quả trình bày trong ĐATN là trung thực và là sản phẩm do tôi tự nghiên cứu, không sao chép từ bất kỳ công trình nào khác. Tất cả các tài liệu tham khảo trong ĐATN – bao gồm hình ảnh, bảng biểu, số liệu và trích dẫn – đều được trích dẫn rõ ràng và đầy đủ trong danh mục tài liệu tham khảo. Tôi hoàn toàn chịu trách nhiệm trước nhà trường nếu có bất kỳ hành vi sao chép nào vi phạm quy định.

\begin{flushright}
    \begin{minipage}{7cm}
        \centering
        Hà Nội, ngày ... tháng ... năm 2025 \\
        \vspace{0.5cm}
        \textbf{Tác giả ĐATN} \\
        \vspace{1.2cm}
        Phạm Quốc Dũng
    \end{minipage}
\end{flushright}
\clearpage
% 3. Tóm tắt nội dung
\chapter*{TÓM TẮT NỘI DUNG ĐỒ ÁN}
\addcontentsline{toc}{chapter}{TÓM TẮT NỘI DUNG ĐỒ ÁN}

Đồ án trình bày quá trình thiết kế và triển khai hệ thống web tự động chấm điểm phiếu trắc nghiệm (OMR) dựa trên ảnh chụp từ điện thoại thông minh, không yêu cầu máy quét chuyên dụng.

\textbf{Bài toán:} Các cơ sở giáo dục hiện nay phụ thuộc vào chấm tay hoặc máy OMR đắt tiền. Hệ thống đề xuất cho phép giáo viên chụp ảnh phiếu bằng điện thoại và nhận kết quả chấm trong vòng 1--3 giây, với chi phí triển khai gần bằng không.

\textbf{Giải pháp kỹ thuật:} Hệ thống được xây dựng theo kiến trúc ba tầng với backend FastAPI (Python), frontend React và cơ sở dữ liệu PostgreSQL. Pipeline xử lý ảnh gồm 15 bước sử dụng OpenCV và NumPy: chuẩn hóa phối cảnh (Perspective Warp), nhị phân hóa thích nghi (Otsu/Adaptive Gaussian), phát hiện marker định vị, suy luận vùng quan tâm (ROI) và giải mã đáp án. Ba đóng góp kỹ thuật nổi bật gồm: (1)~cơ chế \textit{MCQ Map Search Rescue} tự động hiệu chỉnh lưới câu hỏi khi phát hiện lệch lưới bằng cách thử 54 biến thể hình học; (2)~thuật toán tính điểm bong bóng tổng hợp kết hợp mật độ pixel nhị phân và độ tối từ ảnh xám gốc, tăng độ chính xác với bong bóng tô mờ; (3)~tính năng \textit{Smart Camera Scanner} phát hiện phiếu theo thời gian thực trên trình duyệt bằng WebRTC, tự động chụp khi đủ bốn marker và điều kiện ánh sáng.

\textbf{Kết quả:} Hệ thống hỗ trợ đầy đủ luồng nghiệp vụ: tạo bài thi với nhiều mã đề, chấm đơn và chấm lô tối đa 50 ảnh, xem lịch sử và xuất kết quả dạng Excel/PDF. Kiến trúc Form Profile cho phép mở rộng sang nhiều mẫu phiếu khác nhau mà không cần sửa pipeline chính.

\clearpage
% 4. Abstract (English)
\chapter*{ABSTRACT}
\addcontentsline{toc}{chapter}{ABSTRACT}

This thesis presents the design and implementation of a web-based system for automatically grading multiple-choice answer sheets (OMR) from smartphone photos, requiring no dedicated scanner hardware.

\textbf{Problem:} Educational institutions currently rely on manual grading or expensive OMR machines. The proposed system enables teachers to photograph answer sheets with a smartphone and receive grading results within 1--3 seconds, with near-zero deployment cost.

\textbf{Technical Solution:} The system follows a three-tier architecture with a FastAPI (Python) backend, React frontend, and PostgreSQL database. The image processing pipeline consists of 15 steps using OpenCV and NumPy: perspective warping, adaptive binarization (Otsu/Adaptive Gaussian), positioning marker detection, region-of-interest (ROI) inference, and answer decoding. Three notable technical contributions include: (1)~the \textit{MCQ Map Search Rescue} mechanism that automatically corrects the question grid when misalignment is detected by testing 54 geometric variants; (2)~a composite bubble-scoring algorithm combining binary pixel density and darkness from the original grayscale image, improving accuracy for lightly-filled bubbles; (3)~the \textit{Smart Camera Scanner} feature that detects answer sheets in real time via WebRTC in the browser, automatically capturing when all four markers and lighting conditions are satisfied.

\textbf{Results:} The system supports a complete business workflow: creating exams with multiple test codes, single and batch grading of up to 50 images, history viewing, and result export in Excel/PDF format. The Form Profile architecture allows extension to diverse answer sheet templates without modifying the core pipeline.

\textbf{Keywords:} OMR, multiple-choice grading, computer vision, OpenCV, FastAPI, React, WebRTC.

\clearpage
% 5. Mục lục và Danh mục
\tableofcontents
\newpage

\listoffigures
\addcontentsline{toc}{chapter}{DANH MỤC HÌNH VẼ}
\newpage

\listoftables
\addcontentsline{toc}{chapter}{DANH MỤC BẢNG BIỂU}
\newpage

DANH MỤC THUẬT NGỮ VÀ TỪ VIẾT TẮT
\addcontentsline{toc}{chapter}{DANH MỤC THUẬT NGỮ VÀ TỪ VIẾT TẮT}

\renewcommand{\arraystretch}{1.25}
\begin{longtable}{|L{3.2cm}|L{11.2cm}|}
\hline
\textbf{Thuật ngữ} & \textbf{Ý nghĩa} \\
\hline
\endfirsthead

\hline
\textbf{Thuật ngữ} & \textbf{Ý nghĩa} \\
\hline
\endhead
ASGI & Chuẩn giao tiếp bất đồng bộ giữa web server và ứng dụng Python web. ASGI là viết tắt của \textit{Asynchronous Server Gateway Interface}. \\
\hline

OMR & Nhận diện dấu quang học, dùng để phát hiện các ô được tô trên phiếu trả lời trắc nghiệm. OMR là viết tắt của \textit{Optical Mark Recognition}. \\
\hline

MCQ & Câu hỏi trắc nghiệm nhiều lựa chọn, thường gồm các đáp án A, B, C, D. MCQ là viết tắt của \textit{Multiple Choice Question}. \\
\hline

ROI & Vùng quan tâm trên ảnh cần được xử lý, ví dụ vùng MSSV, vùng mã đề, vùng câu hỏi trắc nghiệm hoặc vùng họ tên. ROI là viết tắt của \textit{Region of Interest}. \\
\hline

Marker & Dấu định vị màu đen được in trên phiếu, dùng để xác định vị trí phiếu, căn chỉnh phối cảnh và suy luận các vùng cần đọc. \\
\hline

Perspective Warp & Phép biến đổi phối cảnh, dùng để đưa ảnh phiếu bị chụp nghiêng về mặt phẳng chuẩn. \\
\hline

Binarization & Quá trình nhị phân hóa ảnh, chuyển ảnh xám thành ảnh chỉ gồm hai nhóm chính là nền và vùng tối. \\
\hline

Threshold & Ngưỡng phân tách pixel sáng và pixel tối trong quá trình nhị phân hóa ảnh. \\
\hline

Otsu Threshold & Phương pháp tự động chọn ngưỡng nhị phân hóa dựa trên phân bố mức xám của ảnh. \\
\hline

Adaptive Threshold & Phương pháp nhị phân hóa thích nghi theo từng vùng ảnh, phù hợp khi ảnh có ánh sáng không đều hoặc bóng đổ. \\
\hline

MCQ Map Search Rescue & Cơ chế thử các biến thể nhỏ của lưới câu hỏi để cải thiện kết quả khi số câu không chắc chắn vượt ngưỡng. \\
\hline

Smart Camera Scanner & Tính năng quét phiếu bằng camera trên trình duyệt, tự động phát hiện marker và chụp ảnh khi phiếu đủ điều kiện. \\
\hline

WebRTC MediaDevices API & API của trình duyệt dùng để truy cập camera hoặc microphone. Trong hệ thống này, API được dùng để lấy luồng video camera cho Smart Camera Scanner. \\
\hline

\end{longtable}
\normalsize
\newpage

% ------------------------------------------
% NỘI DUNG CHÍNH
% ------------------------------------------
\pagenumbering{arabic}
